\section{Descrizione e obiettivi}
Lo scopo dell'esperienza è la verifica di alcune leggi di ottica geometrica attraverso un banco ottico, una sorgente luminosa, alcuni 
corpi  ottici, alcune lenti montate su un supporto  fissabile  sul  banco  ottico,  una  piattaforma  circolare  graduata  anch'essa 
fissabile  sul  banco  ottico  e  uno  schermo millimetrato posto in fondo al banco. La sorgente luminosa può fornire uno o cinque fasci di luce bianca o colorata e un'immagine graduata. La luce, che è radiazione elettromagnetica, viene in questo caso trattata come raggio, poiché si propaga attraverso fenditure molto più grandi della sua lunghezza d'onda, rendendo trascurabili i fenomeni di diffrazione. I raggi si propagano in linea retta nei mezzi omogenei e isotropi, ma vengono riflessi o rifratti a seconda della superficie di separazione tra due mezzi che incontrano. \\
La prima parte dell'esperienza consiste nella verifica della legge di Snell (o legge di rifrazione), la reversibilità del cammino ottico e nella misura, attraverso due procedure diverse, l'indice di rifrazione di prismi in plexiglass e vetro. La seconda parte dell'esperienza si propone di verificare la legge dei punti coniugati per lenti sottili in plexiglass e di misurarne l'ingrandimento trasversale.




